\section{CoCos and preferred shares: law beyond the prospectus}
\label{sec:coco-legal-risk}

Suppose Ms Lee now asks to buy a bank capital instrument. Her wealth and her
consent to the streamlined approach do not answer the first product question:
what, exactly, could happen to her investment if the issuing bank becomes
distressed? The answer may depend on the offering conditions, banking legislation,
a regulator's order and subsequent litigation. A second question concerns the
intermediary: which products may it sell to this client, through this service,
and subject to which obligations? Keeping these questions separate is necessary
even when both eventually depend on the same loss-absorption provision.

The preserved prospectuses make this a concrete investigation. The UBS Singapore
dollar notes issued in 2024 contain conversion provisions; the US preferred-share
documents describe different legal interests and distribution rights. A country
label or the word ``preferred'' cannot substitute for identifying the issuer,
instrument, applicable conditions and distribution arrangement. The Swiss and
Hong Kong sources examined below form a bounded survey, inspected on
29 September 2026. A dated source establishes what that source says; it does not
establish that no later amendment, order or judgment exists.

\subsection{What can disappear when a bank absorbs losses?}

A contingent convertible instrument, commonly called a \emph{CoCo}, changes the
holder's rights when specified conditions occur. Bank capital instruments may
instead provide for a reduction of the debt claim. \emph{Additional Tier 1}
(AT1) is a category of regulatory capital available to absorb losses. It is not
a promise of safety to its purchaser. \emph{Common Equity Tier 1} (CET1) is the
core equity category; a contractual trigger may refer to its ratio to
risk-weighted assets. The relevant ratio, entity, threshold, observation date
and notification provisions must all be identified. A \emph{point of
non-viability} provision concerns intervention to prevent failure and need not
coincide with breach of that numerical threshold
\citep[arts.~27, 27a and 29]{swisscao2025}.

A permanent write-down can extinguish principal without delivering replacement
shares. Conversion also removes the original debt claim, but replaces it with
equity under specified settlement arrangements. The market value of that equity
may be very low or eventually zero; this does not make conversion the same legal
mechanism as cancellation without compensation. Temporary reduction and a
conditional future write-up form a third possibility. Separately, cancellation
of a non-cumulative distribution removes that period's payment entitlement
without, by itself, cancelling principal. The implementation must therefore
record debt principal, replacement shares, distribution entitlements and any
restoration right separately. Its existing exact-arithmetic examples do this
conditionally; they do not establish that every real instrument fits those
examples.

The retained UBS conditions illustrate why the qualification matters. Conditions
7 and 8 provide for conversion following the specified trigger or viability
event. Condition 7(b) also contains notification, timing and higher-trigger
priority provisions, including a defined circumstance in which FINMA can agree
that conversion need not occur. Condition 7(e) addresses subsequent legal changes
that can displace the contractual viability clause after the required
determination and notice. Thus, even this conversion example needs more than a
comparison of one capital ratio with one number. The retained document contains
published terms and conditions; it is not represented here as a complete final
prospectus with every incorporated document and later notice
\citep[conditions~7--8, PDF pp.~25--29]{ubsat12024}.

\emph{Bail-in} introduces another route: an authority restructures liabilities
under statutory powers. Whether that route reaches a particular claim depends
on the applicable law and proceedings. A warning about possible bail-in is
therefore evidence of a risk requiring investigation. It is neither evidence
that an order has been made nor a sufficient classification rule for every
investor-protection regime. Hong Kong's banking regulator expressly distinguishes
debt that could be bailed in from the narrower product class covered by its
loss-absorption guidance \citep[questions~1--4]{hkmalacfaq2022}.

\subsection{Swiss capital requirements and powers of intervention}

The Swiss Financial Market Supervisory Authority, \emph{FINMA}, supervises the
relevant banks. Three different sources of loss absorption must be examined:
the capital instrument's agreed conditions, ordinary banking resolution law and
any exceptional measure applicable to the event. They do not supply
interchangeable legal grounds.

The retained German consolidation of the Capital Adequacy Ordinance, effective
1 January 2025, is the principal legislative text used for the following
description; its official-site English translation is retained as a reading
aid. Article 27 distinguishes capital eligibility from the trigger written
into a debt instrument. Article 29 requires provision for loss absorption at
non-viability, through complete reduction of claims or conversion, and specifies
FINMA's role. An instrument may consequently qualify as capital for the bank
because it exposes its holder to loss. This prudential eligibility is a different
question from whether an intermediary may distribute it to Ms Lee
\citep[arts.~27--29]{swisscao2025}.

Article 27a, introduced with effect from 1 January 2025, is especially relevant
to executable interpretation. Once the specified trigger for eligible AT1
liabilities occurs, FINMA orders the reduction or conversion, its timing and
the amount affected. The provision also requires preservation of the
contractual issue or loan conditions and their resulting priority. A program
that extracts only ``FINMA orders the reduction'' loses those conditions.
Conversely, a program evaluating a March 2023 event cannot silently use this
later provision. The retained 2025 consolidation also amends article 29(2);
the earlier wording quoted in the Credit Suisse judgment must remain available
for the historical question
\citep[arts.~27a and 29, amendment footnotes]{swisscao2025}.

Ordinary resolution follows a separate statutory route. Article 25 of the
Banking Act states the conditions
for intervention, including serious liquidity problems, suspected
over-indebtedness and failure to meet capital requirements after the specified
deadline. Articles 26 and 28 distinguish protective measures from restructuring
proceedings. The amendment published as AS~2022~732, effective 1 January 2023,
sets out capital measures in article 30b: its provisions address conversion and
reduction of claims, protected claims, prior absorption of capital, creditor
ordering and specified exceptions. Article 30c governs the restructuring plan,
including creditor protections, and article 31 governs approval, including its
specified compensation provision for systemically important banks
\citep[arts.~30b--31 and commencement provision]{swissbankamend2022}. These
provisions are also visible in the retained 2024 consolidation. A general
instruction power is not interchangeable with this structured resolution route
\citep[arts.~25--31]{swissbank2024}.

Subordinate procedural rules also change. FINMA's Insolvency Ordinance of
20 August 2025 entered into force on 1 October 2025 and repealed the 2012
Banking Insolvency Ordinance. Its transitional provision applies the new
ordinance to proceedings already pending at commencement. The investigation
must therefore inspect transitional rules rather than assume that the law
applicable throughout a proceeding is fixed by its opening date
\citep[arts.~1--2 and 48--50]{swissinsolvency2025}.

There was also an exceptional measure in March 2023. The amendment published
as AS~2023~136 inserted article 5a into the emergency liquidity ordinance.
It addressed FINMA's ability to order the borrower and its financial group to
write down additional Tier 1 capital at the specified credit-approval point.
Its text did not declare all CoCos issued by every Swiss bank worthless.
The amendment entered into force on 19 March 2023 at 20:00
\citep[art.~5a and commencement provision]{swissemergency2023}. Another amendment,
adopted on 6 September and published as AS~2023~495, repealed articles 5--7
with effect from \emph{15 September 2023}. Confusing adoption with commencement
would give the program the wrong boundary. Repeal of a power for new action
also does not, by itself, decide the consequences of an earlier order
\citep[repeals and commencement provision]{swissrepeal2023}.

These provisions establish why a legal investigation is necessary. They do not
justify a universal rule that ``Swiss bank'' and ``AT1'' imply an immediate
write-down to zero. The program needs the legal entity and claim covered by a
power, its effective version, its factual conditions, the required procedure
and any operative order. It must also distinguish the existence of an asserted
legal basis from a final judicial determination of its validity.

\subsection{Credit Suisse: the write-down and its judicial review}

The Credit Suisse case demonstrates the cost of collapsing those distinctions.
FINMA's statement of 23 March 2023 described its instruction to write down
Credit Suisse's AT1 instruments. It relied on the contractual viability event
and the emergency ordinance, and said the Tier 2 instruments were unaffected.
That is a preserved account of the regulator's position and action. It is not
a judgment resolving all objections to their legal basis
\citep{finmaat12023}.

In its partial judgment of 1 October 2025, B-2334/2023, the Swiss Federal
Administrative Court annulled the challenged order. Its reasoning examined the
contractual provisions and the alleged public-law grounds separately. In
particular, the court distinguished assistance addressing liquidity from the
capital-adequacy conditions in the contracts before it. It did not treat the
existence of extraordinary public support as sufficient to erase the rest of
the contractual wording. Its reasoning also distinguished a guarantee benefiting
the acquiring entity from assistance to the entity specified in the contract
\citep[reasons~5.3--5.5, especially 5.4.4--5.4.8]{csat1judgment2025}.

On the public-law grounds, the court found that the protective-measure provision
in article 26 of the Banking Act and the general remedial provision in article
31 of the Financial Market Supervision Act did not provide the required basis
for the disputed extinguishment of claims. It also rejected reliance on the
emergency provision, addressing legal certainty, the constitutional conditions
for emergency action, delegation and property protection. These are the
findings of that court in this partial judgment. They should be preserved with
their procedural status, rather than rewritten as either an uncontested
description of FINMA's powers or a general abolition of bank resolution
\citep[reasons~6.8--6.10, 7.6--7.11, 8 and 10]{csat1judgment2025}.

The subsequent chronology is equally material. FINMA announced an appeal on
15 October 2025. The Administrative Court's announcement of 22 October said
that other pending AT1 cases were suspended and that reversal of the write-down
had not yet been decided \citep{csat1appeal2025,csat1suspension2025}.
A later Federal Supreme Court judgment,
2C\_565/2025 of 5 March 2026, records the separate merits appeals
2C\_661/2025 and 2C\_662/2025 and the grant of suspensive effect on
8 December 2025. But 2C\_565/2025 itself concerns a procedural complaint
about access to submissions and delay; its inadmissibility decision does
\emph{not} dispose of those merits appeals
\citep[facts~D, reasons~1.1--1.4 and operative order]{csat1procedure2026}.

The dossier therefore supports a more precise historical account than either
``the regulator may always erase the bonds'' or ``the investors have been
repaid''. It records the write-down, the partial annulment, the separate appeals
and the reported stay. The latest judicial text inspected for that procedural
history is dated 5 March 2026. A fresh operational determination requires a
check of the merits dockets and any later relief; a search that finds no later
decision is not proof that none exists. Payment actually received is a further
fact, distinct from annulment, an entitlement to relief or a stay.

For testing, this case supplies source documents and changing propositions,
not a human-labelled answer with which the program must agree. A court's
judgment is part of the law and procedural history being represented. Whether
the software has represented its content, scope and consequences correctly
remains a separate obligation. The project does not use human ratings,
annotator consensus or agreement with a model's preferred legal interpretation
as a quality criterion.

\subsection{Which Hong Kong requirement follows from the risk?}

The SFC's circular of 7 December 2018 expressly discusses bonds whose
non-viability loss absorption can follow a capital-ratio event or government
or regulatory action. It connects their complexity to product due diligence,
suitability, disclosure and the client's ability to understand and bear the
risks. This supports investigating statutory and regulatory events alongside
contractual triggers. It does not make ``can lose everything'' a universal
SFC prohibition on sale \citep[discussion of NVLA bonds and selling practices]{sfcnvla2018}.

For a bank acting as a \emph{registered institution}, another source is decisive:
the HKMA's circular of 21 October 2022 and its two annexes. The circular
consolidates and supersedes the earlier 2018 and 2019 guidance
\citep{hkmalac2022}. Annex~1(A)
restricts sale and distribution of in-scope loss-absorption products to
professional investors in both primary and secondary markets. Thus the channel
and intermediary type matter: an SFC circular addressed broadly to
intermediaries and an HKMA requirement for registered institutions must not be
treated as identical statements with identical scope
\citep[applicability and section~A]{hkmalacannex2022}.

The scope must be established before applying that restriction. Annex~1 covers
the specified contingent write-down or conversion features, and certain
products investing mainly in, or closely linked to, such instruments
\citep[applicability]{hkmalacannex2022}.
Annex~2, question~1, includes qualifying AT1, Tier 2 and specified
loss-absorbing-capacity debt under the relevant Hong Kong or equivalent foreign
regimes. It expressly excludes instruments legally constituted as equity,
including preferred shares, and deposits. Questions~2--4 explain why possible
bail-in alone is insufficient. Consequently, an ordinary debt warning does not
automatically place the instrument in scope, and exclusion of an equity-form
preferred share from this particular circular does not establish that the
share is non-complex, suitable, or freely distributable. A depositary share
also requires examination of its own legal interest and underlying security;
its marketing name is insufficient
\citep[questions~1--4]{hkmalacfaq2022}.

Exceptions have their own boundaries. Annex~1(E) removes the specified enhanced
suitability and disclosure measures in sections C and D for loss-absorption
funds while preserving the applicable SFC requirements; it does not remove
section A's professional-investor restriction. Section F grants specified
institutional and qualifying corporate professional investors relief from
sections B, C and D. These are limited forms of relief, not a general
permission for retail distribution
\citep[sections~E--F]{hkmalacannex2022}.

The execution arrangement can also change which intermediary must apply the
circular. Annex~2, question~9, describes a specific arrangement involving an
appropriately regulated adviser or asset manager, an execution broker with
the restricted service and contact role, a written allocation of responsibility
and written notification to the customer. The clarification requires all of
those conditions and does not cover a customer asking that broker directly to
purchase the product. Question~10 separately addresses discretionary portfolio
management on a holistic basis. The procedure must therefore preserve service
roles and conditions rather than attach an unconditional ``execution-only''
exemption to the account \citep[questions~9--10]{hkmalacfaq2022}.

SPI streamlining is a further decision. The SFC's SPI FAQ identifies
loss-absorption debt and related products as requiring a specific product
category, with regard to their characteristics and existing product-specific
requirements. Meeting the SPI wealth threshold alone establishes neither
category knowledge nor satisfaction of a separate selling restriction.
The implementation should report product classification, applicable audience
restrictions, remaining selling obligations and availability of streamlining
as separate conclusions. A client's own mandate may add a restriction, but
that restriction must come from the mandate rather than a presumed preference
attached to the client's wealth
\citep[questions~3--4]{sfcspiprospectusfaq}.

\subsection{A source procedure that can detect the missing law}

An instrument investigation begins with identity: the issuing legal entity,
any relevant guarantor or resolution entity, identifier, issue and amendment
dates, legal form, governing law and intended distribution arrangement. The
document set should include final conditions, the correctly associated base
prospectus, incorporated material, supplements and relevant notices. A later
base prospectus must not be attached to an earlier issue merely because it is
from the same bank. A missing final document leaves a recorded gap.

The next step follows the prospectus's legal references and independently
examines the issuer-jurisdiction legislation and supervisory sources. The
investigation records the provision, its effective interval, the entities and
claims it covers, its conditions and the legal consequence. Ordinary resolution,
contractual non-viability and emergency action remain separate entries. It then
checks relevant orders, proceedings, stays and subsequent judgments. For the
Swiss example, this means identifying the relevant ordinance edition and
checking the Credit Suisse litigation without transferring its contractual
findings to a different UBS issue.

A consolidation date is also insufficient evidence that every amendment has
been incorporated. The archived Banking Act copy labelled 1 January 2023 lacks
the article 30b and 30c headings present in the promulgated amendment effective
on that date and in the retained 2024 consolidation. The discrepancy is retained
as a source-version finding; the account above uses the amendment and the
corroborating consolidation. A mechanical article-presence check detects this
omission, although presence alone cannot establish faithful incorporation of
the provision's meaning.

Distribution law is investigated separately for the place, intermediary,
service and customer. For the present case, that search includes the SFC
complex-product and SPI material and the HKMA circular with both annexes,
including their exceptions and supersession history. Product-specific offering
restrictions and client mandates remain additional dependencies. This is a
bounded source procedure: cross-border recognition, other jurisdictions,
incorporated definitions or a newly discovered restriction can enlarge the
required set. The program must then reopen the affected conclusion rather than
assert that the original list was exhaustive.

Two dates are retained for every changing proposition: when it applies and
when it became available in the decision's evidence history. Downloading a
2025 judgment today must not rewrite what the system knew in March 2023.
Conversely, a present assessment of the 2023 event may need that judgment.
Original bytes, retrieval locations and hashes preserve the evidence used;
page text is a derivative. A changed source invalidates dependent results until
the change is assessed. Unchanged bytes prove neither current law nor complete
retrieval, so an unresolved update check remains visible.

At the end of this procedure, every required proposition is established under
stated premises, refuted, unresolved, or conflicting. For example, the fact
that FINMA has an applicable power cannot supply the missing fact that it has
issued the relevant order. An unresolved contractual interpretation must not
be converted into certainty by a compiler, by several agreeing models, or by
a human verifier. The software can instead state exactly which conditional
conclusion follows and which missing premise prevents an unconditional one.

\subsection{Tests that discriminate between the possible mistakes}

The implemented checks exercise these distinctions with formally stated
premises. They do not use a collection of people's preferred eligibility
answers. For a necessary-condition rule, the independent specification is the
set of complete fact assignments satisfying every condition. Incomplete
evidence describes several possible assignments. If all satisfy the rule, the
conditional result is true; if none do, it is false; otherwise it is unresolved.
Conflicting evidence is reported separately, so an empty set of consistent
assignments cannot produce a vacuous permission.

This gives a precise test of the professional-investor restriction. Under the
bounded four-premise specification, violation occurs when the seller is a
registered institution, the product is in scope, the investor is not a
professional investor and the specific exception is known not to apply.
Enumerating partial assignments tests whether the executable
predicate agrees with that independently stated relation. An unknown exception
is neither assumed available nor silently treated as a known negative fact.
Satisfaction of this single restriction still does not establish permission to
sell under all applicable requirements.

The regression cases are designed around errors that would otherwise survive
a successful compilation:

\begin{longtable}{P{0.30\textwidth}P{0.62\textwidth}}
\toprule
Change in the evidence or case & Required distinction \\
\midrule
\endhead
Same Swiss issuer country, different instrument conditions & Country alone
cannot establish a write-down mechanism or the power's instrument scope. \\
Capital-ratio condition absent, viability route unresolved & Absence of the
numerical trigger cannot establish absence of every loss route. \\
Generic statutory bail-in warning added & The warning cannot establish the
HKMA's narrower direct-debt classification by itself. \\
Preferred share legally constituted as equity & Exclusion from the particular
debt classification does not become permission to sell. \\
One required authority premise removed & The asserted power becomes unresolved
unless another known false premise already defeats it. \\
Emergency provision tested at its start or repeal boundary & Apply the exact
recorded effective interval; adoption and commencement are distinct. \\
Judgment added after the original decision & Preserve historical knowledge;
permit a separate reassessment using later information. \\
Procedural appeal declared inadmissible & Do not mark the separate merits
appeals decided. \\
Annulment recorded, repayment evidence absent & Preserve the annulment and
its finality or stay status without inventing a payment. \\
SPI eligibility changed to true & A failed independent selling requirement
continues to prevent a positive sale conclusion. \\
Applicable law bytes changed or disappeared & Dependent evidence is stale;
the previous conclusion remains historical. \\
\bottomrule
\end{longtable}

The current executable addition checks the authority premises, the bounded
direct-debt classification, the professional-investor restriction, independent
sale conditions, time intervals, proceeding types and source changes. It does
not yet execute the full UBS notice and conversion schedule, all fund and
depositary-share classifications, cross-border resolution recognition, or an
automatic reading of court reasons. Those remain explicit implementation
obligations. The source dossier also records the currentness check still needed
before a live decision; no finite regression suite establishes the absence of
an unknown future legal change.

\subsection{What this adds to independent assurance}

The earlier prospectus campaign preserved documents and checked conditional
mechanics, source quotations and shared-language execution. That work is useful,
but a correctly executed ratio test cannot compensate for a missing statutory
power, a wrong legal entity or an omitted selling restriction. This survey
identifies those dependencies before the translation into executable rules.
The added tests establish properties of their explicit representation, while
the preserved sources expose the remaining interpretation questions.

Generalisation requires the same discipline when the next unfamiliar instrument
arrives. The program should preserve the new sources before adapting its
interpretation, attempt the fixed checks, and report which premises it cannot
establish. If a legal change requires repair, the old result remains associated
with its original sources and method, and the repaired method is tested on a
subsequent untouched case. The defensible promise is conditional correctness
with detectable limits and reproducible revision. Universal legal correctness
on unknown future texts has not been proved.
